\DocumentMetadata{
  lang = pt-BR,
  pdfstandard = A-2b,
  pdfversion = 1.7
}

\documentclass{abntexto-ufc}

\ufcsetup{
  type = doctoral-thesis,
  print-mode = single-sided,
  cover = auto,
  catalog-card = false,
  coat-of-arms = true,
  institution = {Universidade Federal do Ceará},
  center = {Centro de Ciências},
  department = {Departamento de Computação},
  doctoral-graduate-program = {Programa de Pós-Graduação em Ciência da Computação},
  doctoral-degree-field = {Ciência da Computação},
  doctoral-concentration = {Computação Gráfica},
  author = {Nome Completo do Autor},
  title = {Exemplo de tese de doutorado com o abntexto-ufc V3},
  subtitle = {Demonstração da estrutura acadêmica e editorial},
  location = {Fortaleza},
  year = {2026},
  approval-date = {7 de setembro de 2026},
  advisor = {Prof. Dr. Nome do Orientador},
  advisor-institution = {Universidade Federal do Ceará (UFC)},
  advisor-unit = {Departamento de Computação},
  examiner-2 = {Profa. Dra. Nome do Segundo Membro},
  examiner-2-unit = {Departamento de Computação},
  examiner-2-institution = {Universidade Federal do Ceará (UFC)},
  examiner-3 = {Prof. Dr. Nome do Terceiro Membro},
  examiner-3-unit = {Programa de Pós-Graduação em Ciência da Computação},
  examiner-3-institution = {Universidade Federal do Ceará (UFC)},
  examiner-4 = {Profa. Dra. Nome do Quarto Membro},
  examiner-4-unit = {Departamento de Informática},
  examiner-4-institution = {Instituição Externa de Teste (IET)}
}

\begin{filecontents*}{generated-thesis-summary.tex}
Este exemplo demonstra a apresentação de uma tese de doutorado no perfil V3 da classe
\texttt{abntexto-ufc}. O conteúdo é deliberadamente sintético: seu objetivo é permitir a
inspeção visual da capa, folha de rosto, folha de aprovação, resumo, sumário, estrutura
textual, paginação e elementos pós-textuais sem representar uma tese científica real.
\ufcSummaryKeywords{tese; LaTeX; UFC; normalização acadêmica}
\end{filecontents*}

\begin{filecontents*}{generated-thesis-abstract.tex}
This example demonstrates the doctoral-thesis presentation provided by the V3 profile of
\texttt{abntexto-ufc}. The content is intentionally synthetic and is intended for visual
inspection of the academic structure rather than as a scientific thesis.
\keywords{doctoral thesis; LaTeX; UFC; academic standards}
\end{filecontents*}

\begin{document}

\pretextual
\ufcPrintCover
\ufcPrintTitlePage
\ufcPrintApprovalPage
\ufcPrintSummary{generated-thesis-summary}
\ufcPrintAbstract{generated-thesis-abstract}
\ufcPrintTableOfContents

\textual

\section{Introdução}
Uma tese de doutorado apresenta uma contribuição científica original e deve comunicar de
forma clara o problema investigado, a fundamentação, a metodologia, os resultados e as
conclusões. Neste arquivo, o texto é apenas demonstrativo e serve para visualizar o perfil
de tese da classe.

\section{Fundamentação teórica}
Esta seção ilustra a hierarquia textual adotada pelo modelo. O conteúdo real de uma tese
deverá ser substituído pela revisão crítica da literatura pertinente ao tema da pesquisa.

\subsection{Conceitos fundamentais}
Subseções podem ser utilizadas para organizar conceitos, métodos e trabalhos relacionados,
preservando a hierarquia de seções definida pelo perfil institucional.

\section{Metodologia}
A metodologia deve descrever materiais, procedimentos, critérios de avaliação, dados,
experimentos e demais elementos necessários à reprodutibilidade do estudo.

\section{Resultados e discussão}
Os resultados devem ser apresentados e discutidos em relação aos objetivos e à literatura.
Este parágrafo existe somente para produzir uma página textual representativa no exemplo.

\section{Conclusão}
A conclusão sintetiza as contribuições, limitações e perspectivas de continuidade da pesquisa.

\appendix{Exemplo de material elaborado pelo autor}
Este apêndice demonstra um elemento pós-textual produzido pelo próprio autor do trabalho.

\annex{Exemplo de material externo}
Este anexo demonstra a posição de material externo na estrutura do documento.\\
\textbf{Fonte:} Instituição Externa de Teste (2026).

\end{document}
